%% Use-case specifications, one per use case, in ID order. Included by UseCase-Specification.tex.

%% ------------------------------------------------------------------ UC-01
\ucspec{UC-01}{Sign up and onboard}{F1 Account and onboarding}{A. Account and access}{Visitor}{Google identity
(A4)}{High}{B0 sign-up; B1 onboarding}{Phan Bình Nguyên Lâm}

\ucpart{Brief description}
A visitor creates an account, accepts the privacy terms and tells OmniCourse which program and cohort they follow and
what they want from a plan. Every other student use case starts after this one: the degree audit needs the program and
cohort, and the planner starts from the planning goal.

\ucpart{Preconditions}
\begin{itemize}
  \item The visitor is not signed in.
  \item At least one program and cohort exist in the curriculum data (UC-11).
\end{itemize}

\ucpart{Basic flow}
\begin{enumerate}
  \item The visitor opens OmniCourse and chooses \st{Sign up}.
  \item The system asks for an email, a password, a display name, and consent to the terms and the privacy notice. The
        notice states what is stored, for how long, and how to export or delete the account.
  \item The visitor fills in the form, ticks consent and presses \st{Create account}.
  \item The system checks the input: a valid email, a password of 10 to 128 characters, a display name of 1 to 100
        characters, and consent given. It shows the password strength as the visitor types.
  \item The system creates the account, stores the password only as an argon2id hash, records the time of consent and
        signs the visitor in on this device. The visitor is now a student.
  \item The system starts onboarding and asks for the \st{program} and \st{cohort}, offering only those that exist in the
        curriculum data.
  \item The student chooses them.
  \item The system asks for the default planning goal: target graduation term, credits per term and internship term.
        Every field is optional and can be changed later.
  \item The student fills in what they know and presses \st{Next}.
  \item The system offers to import a transcript now (UC-02) or later. The student chooses \st{Later}.
  \item The system saves the answers and opens the dashboard with the degree audit (UC-03).
\end{enumerate}

\ucpart{Alternative flows}
\begin{altflows}
  \item[A1. Email already registered.] At step 5 the email belongs to an existing account. The system says so and offers
        \st{Sign in} (UC-13). No account is created.
  \item[A2. Invalid input.] At step 4 a field fails its check. The system marks each failing field and says what is wrong,
        in words as well as by colour. Nothing is stored. The flow returns to step 3.
  \item[A3. Too many attempts.] At step 3, if one network address sends more than 10 sign-up requests in a minute, the
        system refuses further attempts for the rest of the minute and says so.
  \item[A4. Sign up with Google] (after B0). At step 1 the visitor chooses \st{Continue with Google}. Google confirms the
        identity and a verified email. The system asks for consent as in step 2, creates the account with no password and
        continues at step 6.
  \item[A5. Import the transcript now.] At step 10 the student chooses \st{Import transcript}. UC-02 runs; when it ends,
        the flow continues at step 11.
  \item[A6. Program not offered yet.] At step 7 the student's program or cohort is not listed. The system explains that only
        the listed programs are supported in this release. The student may leave onboarding unfinished; guides (UC-10)
        still work, while the audit and the planner wait until a program is chosen.
  \item[A7. Onboarding interrupted.] The student closes the page between steps 6 and 10. The answers already given are
        kept, and onboarding resumes at the first unanswered question on the next sign-in.
\end{altflows}

\ucpart{Postconditions}
\begin{itemize}
  \item \textbf{Success:} an account exists with a consent time; the student is signed in on this device; program, cohort
        and the default planning goal are stored.
  \item \textbf{Failure:} no account is created (A1 to A3), or the account exists with onboarding unfinished (A6, A7).
\end{itemize}

\ucpart{Special requirements}
\begin{itemize}
  \item Security: argon2id password hashes; sessions rotate, expire after 7 days and can be revoked (UC-14).
  \item Privacy: consent is required and its time stored (Luật 91/2025/QH15). No school-portal password is ever asked for.
  \item Academic rules: onboarding stores the credits per term the student asks for and does not apply the credit cap; the
        rule engine applies it when a plan is built (UC-19).
  \item Usability: works from 360 px wide, by keyboard, in Vietnamese and English.
\end{itemize}

\ucpart{Relationships}
Extended by UC-02 at step 10. Followed by UC-03.

%% ------------------------------------------------------------------ UC-02
\ucspec{UC-02}{Import transcript}{F2 Transcript import}{B. Academic record}{Student}{none}{High}{B2}{Phan Bình Nguyên Lâm}

\ucpart{Brief description}
A student brings their transcript into OmniCourse by uploading the transcript PDF exported from the HCMUS portal, or by
typing the rows in. The student reviews every row before it counts. The saved transcript is what the degree audit, the
grades view and the planner read. A deterministic parser reads the file; no language model is involved (ADR 0003).

\ucpart{Preconditions}
The student is signed in and has chosen a program and cohort (UC-01).

\ucpart{Basic flow}
\begin{enumerate}
  \item The student opens \st{Transcript} and chooses \st{Import from portal PDF}.
  \item The system says which file to download from the portal and how, and that the file is read on the OmniCourse
        server, never sent to an AI or any other service, and deleted as soon as it has been read.
  \item The student drops or chooses the PDF.
  \item The system reads the file and finds one row per course taken: code, name, credits, grade and term.
  \item The system shows a review table with every row. A row whose code is not in the student's curriculum is marked
        \st{Not in your curriculum} and kept, since it may be an extra course.
  \item The student checks the rows and corrects, removes or adds rows (A4).
  \item The student presses \st{Save transcript}.
  \item The system saves the rows as the student's transcript and deletes the uploaded file.
  \item The system opens the degree audit (UC-03), recomputed from the new transcript.
\end{enumerate}

\ucpart{Alternative flows}
\begin{altflows}
  \item[A1. The file cannot be read.] At step 4 the file is not a portal transcript, has no text layer, or has a layout the
        parser does not know. The system says why and offers manual entry (A4). Nothing is stored; the file is deleted.
  \item[A2. A row is incomplete.] At step 4 one field of a row cannot be read. The system shows the row with that field
        marked. At step 7 the system refuses to save until the row is completed or removed.
  \item[A3. Course in progress.] At step 4 a row of the current term has no grade yet. The system marks it
        \st{In progress} and keeps it.
  \item[A4. Manual entry.] At step 1, or from A1, the student chooses \st{Enter manually}; at step 6 they choose
        \st{Add a row}. The system offers the courses of the curriculum to pick from and accepts a code that is not in it.
        Credits come from the curriculum when the course is in it. The flow returns to step 6.
  \item[A5. Re-import over a saved transcript.] At step 5 the student already has a transcript. The system shows the rows
        added, changed and removed. At step 8 the new rows replace the old ones, and UC-12 may tell the student what changed
        in their audit.
  \item[A6. The same course twice.] A course was failed and retaken. The system keeps both rows. Which attempt counts is
        decided by the rule engine (UC-19), not by the import.
  \item[A7. Wrong file type or size.] At step 3 the file is not a PDF or is larger than 5 MB. The system refuses it and
        says what it accepts.
  \item[A8. Cancel.] Before step 7 the student leaves. Nothing is stored and the file is deleted. A transcript saved earlier
        is unchanged.
\end{altflows}

\ucpart{Postconditions}
\begin{itemize}
  \item \textbf{Success:} the student's transcript is the reviewed set of rows, and the uploaded file no longer exists.
  \item \textbf{Failure or cancel:} the earlier transcript, if any, is unchanged; no uploaded file remains.
\end{itemize}

\ucpart{Special requirements}
\begin{itemize}
  \item Determinism: the same file always gives the same rows; the parser is unit-tested on synthetic transcripts only.
  \item Privacy: the file is processed in memory or a temporary file, never logged, never sent to a third party.
  \item Usability: each problem (not in curriculum, incomplete, in progress) is marked in words, not by colour alone.
\end{itemize}

\ucpart{Relationships}
Extends UC-01 at its step 10. Feeds UC-03, UC-15, UC-05 and UC-08.

%% ------------------------------------------------------------------ UC-03
\ucspec{UC-03}{View degree audit}{F3 Degree audit dashboard}{B. Academic record}{Student}{none}{High}{B1}{Võ Nguyễn
Minh Triết}

\ucpart{Brief description}
The student sees exactly where they stand: credits passed, in progress and still needed in each requirement group, the
cumulative GPA, the estimated graduation term and a short list of what needs attention. Every figure comes from the rule
engine, applied to the student's transcript and their cohort's curriculum.

\ucpart{Preconditions}
The student is signed in and has chosen a program and cohort (UC-01). A transcript is optional (A1).

\ucpart{Basic flow}
\begin{enumerate}
  \item The student opens the \st{Dashboard}.
  \item The system asks the rule engine (UC-19) to evaluate the transcript against the curriculum of the student's
        cohort.
  \item The system shows the summary: cumulative GPA with its change from the last term, credits passed out of the
        total required, credits left and the estimated graduation term.
  \item The system shows a chart of credits by requirement group, each group in its own colour and glyph: passed, then in
        progress, then still to take.
  \item The system lists the groups with a bar each: credits passed, in progress and required.
  \item The system shows \st{Needs attention}: for example an elective group that is still empty, a failed course not yet
        retaken, a prerequisite missing for next term, or a language requirement not met. Each item names the rule and
        links to the place to act (planner, guide or transcript).
  \item The student opens one group.
  \item The system lists that group's courses in three parts: passed, in progress and still needed, with the credits of
        each.
  \item The student reads the result and leaves, or opens a course's grades (UC-15).
\end{enumerate}

\ucpart{Alternative flows}
\begin{altflows}
  \item[A1. No transcript yet.] At step 2 the student has no transcript. The system shows every requirement as still
        needed, says the audit assumes no course has been passed, and offers \st{Import transcript} (UC-02).
  \item[A2. Look at a past term.] At step 3 the student picks an earlier term from the term picker or hovers a past term.
        The system shows the whole dashboard as it stood at the end of that term. The arrows or \st{Now} return to the
        present.
  \item[A3. Course outside the curriculum.] At step 2 a passed course belongs to no requirement group. The rule engine
        reports it as not counted, and the system lists it under the summary as \st{Not counted toward graduation}.
  \item[A4. Graduation status reached.] At step 6 every requirement is met. The system says the student meets the
        graduation requirements of the curriculum, cites the regulation, and links the graduation guide (UC-10).
\end{altflows}

\ucpart{Postconditions}
Nothing is changed. The student has seen their status as the rule engine computes it today.

\ucpart{Special requirements}
\begin{itemize}
  \item Performance: the dashboard loads in under 2 seconds for a 40-course transcript.
  \item Accuracy: the audit uses the curriculum version of the student's own cohort; tests check each group against
        synthetic transcripts with known results.
  \item Usability: colour is never the only signal; each group has a glyph and every figure has a text label.
\end{itemize}

\ucpart{Relationships}
Includes UC-19. Extended by UC-15.

%% ------------------------------------------------------------------ UC-04
\ucspec{UC-04}{Explore curriculum}{F4 Curriculum explorer}{B. Academic record}{Student}{none}{High}{B1}{Võ Nguyễn
Minh Triết}

\ucpart{Brief description}
The student browses every course of their curriculum by requirement group, with its credits, prerequisites, what it
unlocks and when it is offered, and sees which courses they can take now. A prerequisite graph shows the chains.

\ucpart{Preconditions}
The student is signed in and has chosen a program and cohort (UC-01).

\ucpart{Basic flow}
\begin{enumerate}
  \item The student opens \st{Curriculum}.
  \item The system asks the rule engine (UC-19) for the state of each course for this student.
  \item The system lists the courses group by group. In each group it shows, in order: in progress, available now,
        locked, and passed. A locked course shows a padlock and the prerequisites still missing.
  \item The student types in the search box, for example \texttt{cau truc du lieu}.
  \item The system narrows the list as they type, matching codes and names with or without Vietnamese diacritics.
  \item The student opens a course.
  \item The system shows the course drawer: code, name, credits, group, prerequisites with their state, the courses it
        unlocks, the terms it is offered, and a link to place it in a plan (UC-05).
  \item The student closes the drawer and continues browsing, or leaves.
\end{enumerate}

\ucpart{Alternative flows}
\begin{altflows}
  \item[A1. Prerequisite graph.] At step 3 the student switches to \st{Graph}. The system draws the courses as nodes and
        the prerequisites as arrows. Clicking a node lights its whole chain, before and after; dragging or scrolling pans
        the graph.
  \item[A2. Nothing matches.] At step 5 no course matches. The system says so and keeps the filter visible so the student
        can clear it.
  \item[A3. Offering unknown.] At step 7 the curriculum data does not give the offering terms of the course. The system
        says the terms are not known yet, instead of guessing.
\end{altflows}

\ucpart{Postconditions}
Nothing is changed.

\ucpart{Special requirements}
\begin{itemize}
  \item Search finds a course from a code, a name or a name typed without diacritics, in Vietnamese or English.
  \item The graph and the list are keyboard operable; the graph has a list alternative for screen readers.
\end{itemize}

\ucpart{Relationships}
Includes UC-19. Links to UC-05.

%% ------------------------------------------------------------------ UC-05
\ucspec{UC-05}{Build and edit a plan}{F5 Semester roadmap planner}{C. Planning}{Student}{none}{High}{B2}{Lê Hữu
Hoà}

\ucpart{Brief description}
The student builds a plan to graduation by placing courses into terms, one row per term. Every change is checked at once
by the rule engine against prerequisites, offering terms, the credit cap and requirement coverage, and the result shows
before the change is made. The student can keep several plans.

\ucpart{Preconditions}
\begin{itemize}
  \item The student is signed in and has chosen a program and cohort (UC-01).
  \item The curriculum of the cohort is available (UC-11).
\end{itemize}

\ucpart{Basic flow}
\begin{enumerate}
  \item The student opens \st{Planner}.
  \item The system shows the active plan: the terms already passed (summarised in one row), the current term, and one row
        per future term up to the target graduation term. Each row shows its courses as tiles, its credit load and a
        meter with the program's credit minimum and cap. A strip at the top shows the estimated graduation term and the
        coverage of every requirement group. The \st{Issues} panel on the right lists the violations.
  \item The student picks a course, from the list of courses still needed or from another term, and starts to drag it.
  \item While the course is dragged, the system asks the rule engine (UC-19) about each term and shows the answer:
        terms that cannot take the course (not offered, or before a prerequisite) fade out and refuse the drop; the
        target row previews its new credit load; every row turns green or red with the reason; the \st{Issues} panel
        previews which issues would appear and which would be resolved.
  \item The student drops the course on a term.
  \item The system places the course, re-evaluates the plan, updates the meters, the coverage strip, the estimated
        graduation term and the \st{Issues} panel, and saves the plan.
  \item The student repeats steps 3 to 6 until the plan suits them, then leaves. The plan stays saved.
\end{enumerate}

\ucpart{Alternative flows}
\begin{altflows}
  \item[A1. Keyboard.] At step 3 the student focuses a tile and presses Space to lift it. The arrow keys move it between
        terms, skipping terms that refuse it, and Enter drops it. Escape cancels.
  \item[A2. Fix suggestion.] At step 6 the plan has a violation. The student opens it in the \st{Issues} panel. The
        system shows the rule, the regulation clause and, when the rule engine finds one, a fix, such as moving a
        course to a later term. The student presses \st{Apply}; the system applies it and returns to step 6.
  \item[A3. Undo.] After step 6 the student presses \st{Undo}. The system restores the plan as it was before the last
        change.
  \item[A4. Remove a course.] At step 3 the student drags a placed course back to the list of courses still needed. The
        system removes it from the term and returns to step 6.
  \item[A5. Over the cap.] At step 5 the drop takes the term above the credit cap. The drop is allowed, the row turns red
        with the reason, and the issue is listed until the student fixes it.
  \item[A6. Several plans.] At step 2 the student opens the plan menu to create, rename, duplicate or switch plans. Deleting
        a plan asks for confirmation first, because it cannot be undone.
  \item[A7. Plan that already breaks a rule.] At step 2 a saved plan breaks a rule, for example after the curriculum was
        corrected (UC-11). The plan still loads and the issue is listed; nothing is moved for the student.
  \item[A8. Start from a goal.] At step 2 the student has no plan yet. The system offers an empty plan or
        \st{Generate a plan} (UC-07).
\end{altflows}

\ucpart{Postconditions}
\begin{itemize}
  \item \textbf{Success:} the plan is saved with every change; its violations, coverage and estimated graduation term are
        those the rule engine computed for the final state.
  \item The transcript and the other plans are unchanged.
\end{itemize}

\ucpart{Special requirements}
\begin{itemize}
  \item Performance: a change is validated in under 300 ms, so the preview in step 4 follows the pointer.
  \item Academic rules: the planner never decides a rule; terms fade, rows turn red and issues appear only because the rule
        engine said so (UC-19).
  \item Usability: drag and drop has a full keyboard path (A1); every red or green state has a text reason.
\end{itemize}

\ucpart{Relationships}
Includes UC-19. Extended by UC-06. Starts UC-09 from the plan menu. Opened by UC-07 after a plan is generated.

%% ------------------------------------------------------------------ UC-06
\ucspec{UC-06}{Simulate a what-if scenario}{F7 What-if simulator}{C. Planning}{Student}{none}{Medium}{B3}{Lê Hữu Hoà}

\ucpart{Brief description}
Before making a choice, the student sees what it would cost: failing a course, changing specialization, adding a summer
term or taking a heavier load. The scenario runs on a copy of the plan, and the system shows the new violations and how the
estimated graduation term moves.

\ucpart{Preconditions}
The student is signed in and has a saved plan open in the planner (UC-05).

\ucpart{Basic flow}
\begin{enumerate}
  \item In the planner, the student presses \st{What if}.
  \item The system makes a copy of the plan and offers the scenarios: \st{Fail a course}, \st{Change specialization},
        \st{Add a summer term} and \st{Take more credits}.
  \item The student chooses \st{Fail a course} and picks a course of the current term.
  \item The system asks the rule engine (UC-19) to evaluate the copy with that course failed: the course must be retaken,
        and every course that needs it moves out of the terms it can no longer be in.
  \item The system shows the copy beside the original: the courses that moved, the new violations, and the estimated
        graduation term before and after, for example \emph{Summer 2028 instead of Spring 2028}.
  \item The student presses \st{Save as a new plan} and gives it a name.
  \item The system saves the copy as a new plan. The original plan is unchanged.
\end{enumerate}

\ucpart{Alternative flows}
\begin{altflows}
  \item[A1. Change specialization.] At step 3 the student picks another specialization. The rule engine re-evaluates
        coverage against that specialization's groups; the system shows the courses that no longer count and those now
        needed.
  \item[A2. Add a summer term.] At step 3 the student adds a summer term. The rule engine checks the summer credit cap and
        which courses are offered in summer; the system shows how much earlier the student could graduate.
  \item[A3. Take more credits.] At step 3 the student raises the credits per term. The rule engine keeps every term within
        the program's cap; the system shows the earliest graduation term the new load allows.
  \item[A4. Combine scenarios.] At step 5 the student adds another scenario to the same copy. The flow returns to step 4.
  \item[A5. Discard.] At step 6 the student presses \st{Discard}. The copy is deleted; nothing is saved.
  \item[A6. Data missing.] At step 4 a scenario needs data the curriculum does not have, such as a specialization that is
        not entered yet. The system says what is missing and does not invent it.
\end{altflows}

\ucpart{Postconditions}
\begin{itemize}
  \item \textbf{Saved:} a new plan exists; the original plan and the transcript are unchanged.
  \item \textbf{Discarded:} nothing is stored.
\end{itemize}

\ucpart{Special requirements}
The scenario never changes the transcript or the original plan. All consequences come from the rule engine.

\ucpart{Relationships}
Extends UC-05. Includes UC-19.

%% ------------------------------------------------------------------ UC-07
\ucspec{UC-07}{Generate a plan}{F6 AI plan generator (goal box)}{C. Planning}{Student}{Model provider (steps 3 to 5
only)}{High}{B3}{Phan Bình Nguyên Lâm}

\ucpart{Brief description}
The student states a planning goal in a form, or in their own words in the goal box, which a language model turns into the
form's fields. A deterministic planner then builds a complete plan to graduation that breaks no academic rule and explains
why each course is where it is. The model never places a course.

\ucpart{Preconditions}
\begin{itemize}
  \item The student is signed in and has chosen a program and cohort (UC-01).
  \item The curriculum and the rule engine are available.
\end{itemize}

\ucpart{Basic flow}
\begin{enumerate}
  \item The student opens \st{Generate a plan}.
  \item The system shows the planning-goal form, filled from the student's default goal: target graduation term, credits
        per term, internship term, courses to take and courses to avoid. Above it, the goal box invites a goal in the
        student's own words; a note says the text goes to an AI provider, so names and student IDs should be left out.
  \item The student types, for example, \emph{graduate by summer 2028, at most 20 credits a term, no networking courses},
        and presses \st{Fill the form}.
  \item The system sends the sentence and the course codes and names of the curriculum to the model through the AI
        gateway, and asks for the form's fields as structured data.
  \item The system checks the reply against the form's schema, fills the form and highlights every field that changed.
  \item The student reviews the form, corrects any field and presses \st{Plan}.
  \item The deterministic planner builds a plan from the next term to the target term, asking the rule engine (UC-19) to
        check every placement. The system shows progress.
  \item The system shows the plan term by term with each term's credits. Each course carries the planner's reason, such as
        \emph{prerequisite CS162 is passed the term before} or \emph{offered in main terms only}. The plan shows its
        estimated graduation term.
  \item The student presses \st{Save plan} and names it.
  \item The system saves it next to the student's other plans and opens it in the planner (UC-05).
\end{enumerate}

\ucpart{Alternative flows}
\begin{altflows}
  \item[A1. Without the goal box.] At step 3 the student edits the form directly. The flow continues at step 6; no model is
        involved.
  \item[A2. AI unavailable.] At step 4 the gateway answers \texttt{ai\_unavailable}: no provider, a timeout, a reply that
        still does not fit the schema after one retry, or the monthly budget spent. The system says the goal box is
        unavailable, keeps the form as it was and continues as A1. Planning never depends on the model.
  \item[A3. Daily AI limit reached.] At step 4 the student has used their daily AI calls. The system says when the goal box
        works again and continues as A1.
  \item[A4. A field the rules forbid.] At step 6 a field asks for more than the rules allow, such as credits above the
        cap. The form keeps the value; at step 7 the planner stays within the cap, and at step 8 the system says which part
        of the goal it could not meet, citing the rule. The model never changes a value to fit a rule.
  \item[A5. Goal cannot be met.] At step 7 no plan reaches the target term. The planner returns the plan with the earliest
        graduation term it can reach, and the system names the rules that make the target impossible, such as a
        prerequisite chain that is too long.
  \item[A6. A course to avoid is compulsory.] At step 7 the planner places it anyway and explains at step 8 that every
        student of the cohort must pass it.
  \item[A7. Planning takes too long.] At step 7 the planner has not finished within 30 seconds. The student can cancel at
        any time while progress shows.
  \item[A8. Discard.] At step 9 the student leaves without saving. No plan is stored.
\end{altflows}

\ucpart{Postconditions}
\begin{itemize}
  \item \textbf{Success:} a new plan is saved, and it broke no academic rule when it was generated.
  \item \textbf{Discard or failure:} no plan is stored; the other plans are unchanged.
  \item Each call that reached the model leaves one usage record (user, feature, model, tokens, time) with no text.
\end{itemize}

\ucpart{Special requirements}
\begin{itemize}
  \item Academic rules: every placement and reason comes from the planner and the rule engine. The model fills form fields
        only, and the student sees them before they have any effect.
  \item Privacy: only the sentence and the curriculum's codes and names go to the model; no name, email, student ID or
        transcript. Prompt and reply content is never logged.
  \item Performance: a plan returns within 30 seconds, with visible progress.
\end{itemize}

\ucpart{Relationships}
Includes UC-19. Hands the saved plan to UC-05.

%% ------------------------------------------------------------------ UC-08
\ucspec{UC-08}{Ask the advisor chat}{F8 AI academic advisor chat}{D. Help and reminders}{Student}{Model
provider}{Medium}{B3}{Nguyễn Anh Dũng}

\ucpart{Brief description}
The student asks a question about their studies in their own words, such as \emph{Can I take Machine Learning before
Probability and Statistics?} The advisor chat answers from the official curriculum and regulations and cites the document
and clause. Anything about the student's own record is checked by the rule engine, and the answer shows that result.

\ucpart{Preconditions}
\begin{itemize}
  \item The student is signed in.
  \item The regulation documents the chat answers from are loaded.
\end{itemize}

\ucpart{Basic flow}
\begin{enumerate}
  \item The student opens \st{Advisor chat}.
  \item On first use only, the system explains what is sent to the AI provider (program, cohort, course codes, grades and
        terms; never name, email or student ID), that the provider does not train on it, and that the content is not
        logged. The student presses \st{I agree}, and the system records the time of consent.
  \item The system shows an empty conversation with example questions and a note that answers can be wrong and the cited
        clause is the authority.
  \item The student types a question and presses \st{Send}.
  \item The system sends the question, the earlier messages and the chat's tools to the model through the gateway.
  \item The model calls the tools it needs, within 5 model turns: search the regulations, look up a course, read the
        student's record, ask the rule engine to check a rule (UC-19), or find a guide. The system shows what is happening,
        for example \emph{Checking prerequisites\ldots}
  \item The system shows the answer as it is written.
  \item When the answer is complete, the system shows each cited clause as a link and, for a statement about the student,
        the rule engine result it came from.
  \item The student asks a follow-up question (step 4) or leaves.
\end{enumerate}

\ucpart{Alternative flows}
\begin{altflows}
  \item[A1. Consent declined.] At step 2 the student presses \st{Not now}. The chat answers general questions about the
        regulations only; the student's record is never sent. A question about their own record gets a reply that it needs
        consent, with a button to give it.
  \item[A2. AI unavailable.] At step 5 the gateway answers \texttt{ai\_unavailable}. The system says the chat is
        unavailable now, keeps the question in the box, and points to the guides (UC-10) and the explorer (UC-04).
  \item[A3. Failure mid-answer.] At step 7 the provider fails. The system keeps the text received, marks it
        \st{Incomplete} and offers \st{Ask again}.
  \item[A4. Daily limit reached.] At step 5 the student has used their daily AI calls (20 by default, shared with UC-07).
        The system says when the chat works again.
  \item[A5. Stop.] At step 7 the student presses \st{Stop}. The answer is marked \st{Stopped}; the call still counts.
  \item[A6. No source found.] At step 6 nothing answers the question. The chat says it found no document and does not
        guess; it suggests the academic office or a matching guide.
  \item[A7. The rule engine cannot decide.] At step 6 the rule or data is not encoded. The chat says it cannot confirm the
        answer for this student, shows the clause if found, and suggests confirming with the academic office.
  \item[A8. A task on a portal.] The question is how to do something on a portal. The chat links the guide (UC-10) instead
        of writing its own steps.
  \item[A9. A request to act.] The student asks the chat to change a plan or register for a course. The chat has no tool
        that changes anything; it explains what the student can do and links the planner (UC-05) or the registration
        guide (UC-10).
  \item[A10. Outside studies.] The question is not about studies. The chat says what it can help with and does not
        answer.
\end{altflows}

\ucpart{Postconditions}
\begin{itemize}
  \item \textbf{Success:} the student has an answer with its sources. Nothing in the transcript, plans or account changed.
  \item Each call that reached the model leaves one usage record without content. Conversations are not stored on the
        server; reloading the page starts a new one.
\end{itemize}

\ucpart{Special requirements}
\begin{itemize}
  \item Citations: every rule stated is cited; links are built only from clauses the search actually returned.
  \item Performance: the first words appear in about one second; each provider request times out after 30 seconds.
  \item Robustness: when the model is down, only the chat and the goal box are unavailable.
\end{itemize}

\ucpart{Relationships}
Includes UC-19. Links to UC-10, UC-05 and UC-02.

%% ------------------------------------------------------------------ UC-09
\ucspec{UC-09}{Share or export a plan}{F9 Plan export and sharing}{C. Planning}{Student}{none}{Low}{B2}{Nguyễn Anh Dũng}

\ucpart{Brief description}
The student gives one of their plans to someone else, usually an advisor: as a PDF file, or as a read-only link that
anyone holding it can open without an account. The student can stop sharing at any time.

\ucpart{Preconditions}
The student is signed in and has at least one saved plan (UC-05 or UC-07).

\ucpart{Basic flow}
\begin{enumerate}
  \item The student opens a plan and presses \st{Share}.
  \item The system shows what a reader will see: the plan's name, program, cohort and curriculum version; each term with
        its courses and credits; the estimated graduation term and the violations; when the plan was last changed. It also
        says what a reader will not see: grades, email, transcript, student ID and other plans. The display name is hidden
        unless the student ticks \st{Show my name}. A warning says anyone with the link can open it.
  \item The student presses \st{Create link}.
  \item The system creates an unguessable link for this plan and shows it with \st{Copy}.
  \item The student copies the link and sends it through any channel. OmniCourse sends nothing itself.
\end{enumerate}

\ucpart{Alternative flows}
\begin{altflows}
  \item[A1. Download as PDF.] At step 1 the student presses \st{Download PDF}. The system makes an A4 PDF with the content
        of step 2 and the date it was made. No link is created and nothing is kept on the server.
  \item[A2. Stop sharing.] The student opens \st{Share} and presses \st{Stop sharing}. The link stops working at once.
  \item[A3. Already shared.] At step 1 the plan has a live link. The system shows it with \st{Copy} and \st{Stop sharing}
        instead of making a second one. A plan has at most one live link.
  \item[A4. Share again.] After A2 the student presses \st{Create link}. A new link is made; the old one stays dead.
  \item[A5. Plan deleted or account deletion requested.] Every link of that plan, or of that student, stops working.
\end{altflows}

\ucpart{Postconditions}
\begin{itemize}
  \item \textbf{Shared:} a live link exists until the student stops sharing or deletes the plan or account.
  \item \textbf{Downloaded:} a PDF is on the student's device; nothing new is stored.
\end{itemize}

\ucpart{Special requirements}
The link holds 32 random bytes, as strong as a session token. The shared page is not indexed by search engines. The PDF is
readable in black and white: violations are marked in words.

\ucpart{Relationships}
Started from the plan menu of UC-05. The advisor's side is UC-16.

%% ------------------------------------------------------------------ UC-10
\ucspec{UC-10}{Follow a school-task guide}{F11 Quicklinks and step-by-step guides}{D. Help and reminders}{Visitor (or
Student)}{HCMUS portal}{Medium}{B1 reading; B3 personal parts}{Nguyễn Anh Dũng}

\ucpart{Brief description}
A visitor or student needs to do a school task that OmniCourse does not perform, such as registering for courses or paying
tuition. They find the guide for it, read when to do it and what to prepare, and follow the steps on the official portal
through the guide's direct link. OmniCourse never signs in to a portal.

\ucpart{Preconditions}
At least one guide is published. No sign-in is needed (ADR 0010).

\ucpart{Basic flow}
\begin{enumerate}
  \item The visitor opens \st{Guides}.
  \item The system lists the guides by task group (Registration, Tuition and fees, Records and certificates, Graduation).
        Each entry shows the task, the portal, the account it needs and the direct link.
  \item The visitor types a few words, such as \texttt{hoc phi}. The system narrows the list as they type, in Vietnamese or
        English, with or without diacritics.
  \item The visitor opens a guide.
  \item The system shows what the task is for, when to do it, what to prepare, the portal and account, numbered steps
        quoting the portal's buttons, common mistakes, and the guide's version and the date it was last checked.
  \item The visitor presses \st{Open} with the portal's name. The system opens the deepest portal page the school allows
        in a new tab.
  \item The visitor signs in on the portal with their school account and completes the task there. OmniCourse sees none of
        it.
\end{enumerate}

\ucpart{Alternative flows}
\begin{altflows}
  \item[A1. Straight to the link.] At step 2 the visitor presses the direct link in the list. The flow continues at step 7.
  \item[A2. No guide matches.] At step 3 nothing matches. The system says so and shows the task groups to browse.
  \item[A3. A value that depends on the program.] At step 5 a note depends on the program, such as the credit cap. A
        signed-in student with a program sees the value from the rule engine (UC-19) with its clause; a visitor is told the
        value depends on the program and is offered to sign in.
  \item[A4. Registration from a plan.] At step 5 of the registration guide, a signed-in student picks a plan. The system
        lists the courses the plan places next term, with codes ready to copy. It never registers for the student.
  \item[A5. Possibly out of date.] At step 5 the guide was last checked more than 6 months ago. The system warns that the
        portal may have changed; the guide stays readable.
  \item[A6. Report a problem.] At any step a signed-in student presses \st{Report a problem}, picks a reason and may add a
        note. The form warns not to include a name, student ID or password. The report is stored against the guide and its
        version.
\end{altflows}

\ucpart{Postconditions}
OmniCourse stores nothing about the task, not even whether it was done. A report made in A6 is stored against the guide.

\ucpart{Special requirements}
No portal access: no portal password asked, no portal form filled, no portal page read. Guides are versioned; every
published change has a new version, a change note and a new last-checked date.

\ucpart{Relationships}
Linked from UC-03 (A4), UC-08 (A2, A8, A9) and UC-12.

%% ------------------------------------------------------------------ UC-11
\ucspec{UC-11}{Maintain curriculum}{F12 Curriculum administration}{E. Curriculum administration}{Program
administrator}{none}{Medium}{B2}{Nguyễn Hữu Quốc Thắng}

\ucpart{Brief description}
A program administrator adds and updates the curriculum of a program they are allowed to maintain, from the official
documents: cohorts, requirement groups with their credit minimums, the courses linked to each group, prerequisites and
offering terms. Each cohort has its own curriculum. No AI model reads or decides the data.

\ucpart{Preconditions}
\begin{itemize}
  \item The administrator is signed in and holds program-administrator permission for the program (UC-18).
  \item They have the official curriculum document for the program and cohort.
\end{itemize}

\ucpart{Basic flow}
\begin{enumerate}
  \item The administrator opens \st{Curriculum administration} and selects a program and cohort. The system offers only
        programs they may maintain.
  \item The system shows the curriculum: groups with their required credits and the courses linked to each.
  \item The administrator chooses an item to add or edit, such as a requirement group, and enters the values from the
        official document. Courses are picked from the shared catalog.
  \item The administrator checks the entries against the document and presses \st{Save}.
  \item The system checks permission, required fields, the course-code format (\texttt{\^{}[A-Z0-9]\{2,20\}\$} after
        trimming and upper-casing) and links to other items, and asks the rule engine (UC-19) for the academic checks, such
        as a prerequisite cycle.
  \item The system saves the whole change at once and shows the saved values, program and cohort.
\end{enumerate}

\ucpart{Alternative flows}
\begin{altflows}
  \item[A1. Add a cohort.] At step 1 the cohort does not exist. The administrator presses \st{Add cohort} and enters it.
        The system rejects a duplicate cohort; earlier cohorts are unchanged.
  \item[A2. Invalid or missing values.] At step 5 a field is empty, a code is duplicated or malformed, or a linked item does
        not exist. The system marks the fields, refuses to save and keeps the saved curriculum unchanged.
  \item[A3. Prerequisite cycle.] At step 5 the rule engine finds that the new prerequisite closes a cycle. The system shows
        the cycle and refuses to save.
  \item[A4. No permission or expired sign-in.] The system refuses the request, says which permission is missing or asks
        the user to sign in again. The backend checks every request, not only the screen.
  \item[A5. Remove an item.] At step 3 the administrator removes a course from a group. The system shows what will be
        removed and asks for confirmation. The course itself and its other links stay. A deletion that other data still
        needs is refused with the list of what depends on it.
  \item[A6. Cancel.] Before step 4 the administrator cancels. The edit is discarded.
  \item[A7. Save fails or the result is unknown.] At step 6 the server fails, or the connection drops. A failed save leaves
        no partial change; if the result is unknown, the administrator reloads the item before trying again.
  \item[A8. The course is not in the catalog.] At step 3 the course does not exist yet. An administrator who also holds
        catalog permission creates it (UC-17) and returns to step 3; otherwise they ask a catalog maintainer.
\end{altflows}

\ucpart{Postconditions}
\begin{itemize}
  \item \textbf{Saved:} the change belongs to the chosen program and cohort and is read by the audit and planner from then
        on. Stored transcripts and saved plans are not rewritten; a plan that now breaks a rule shows the issue (UC-05 A7).
  \item \textbf{Rejected or cancelled:} nothing is saved.
\end{itemize}

\ucpart{Special requirements}
Every write checks permission on the server. Tests use made-up curricula. A second program or cohort entered through this
screen must not change the first (B3 check).

\ucpart{Relationships}
Includes UC-19. Extended by UC-17. Depends on rights given in UC-18.

%% ------------------------------------------------------------------ UC-12
\ucspec{UC-12}{Receive notifications}{F10 Notifications}{D. Help and reminders}{Student}{Time}{Low}{B2}{Nguyễn Anh Dũng}

\ucpart{Brief description}
OmniCourse tells the student, inside the app, about what needs their attention: the registration window opening, a
check-in at the start of each term, a change in their transcript, or a plan that now breaks a rule. The student reads them
in the inbox or the bell popover.

\ucpart{Preconditions}
The student is signed in and has finished onboarding (UC-01).

\ucpart{Basic flow}
\begin{enumerate}
  \item Time reaches a scheduled moment, for example seven days before a registration window opens.
  \item The system asks the rule engine (UC-19) what applies to this student: the courses their active plan places in the
        coming term and any violation it has.
  \item The system creates a notification: \emph{Registration for Spring 2027 opens on 9 December. Your plan places 5
        courses, 19 credits}, with links to the plan and to the registration guide (UC-10).
  \item The student signs in. The bell in the sidebar shows the number of unread notifications.
  \item The student opens the bell. The system shows a short popover of what needs attention and a link to the inbox.
  \item The student opens the inbox. The system lists every notification, newest first, with filters on the right.
  \item The student opens a notification and follows its link. The system marks it read.
\end{enumerate}

\ucpart{Alternative flows}
\begin{altflows}
  \item[A1. Transcript changed.] Instead of step 1, the student saves a new transcript (UC-02 A5). The system creates a
        notification that lists what changed in the audit.
  \item[A2. Plan now breaks a rule.] Instead of step 1, a curriculum change (UC-11) makes the student's active plan break a
        rule. The system creates a notification that names the rule and links the plan.
  \item[A3. Mark all read.] At step 6 the student presses \st{Mark all as read}.
  \item[A4. Filter.] At step 6 the student picks a filter (unread, registration, transcript, plan). The list narrows.
  \item[A5. Turn a kind off.] In settings (UC-14), the student turns off a kind of notification. No new notification of
        that kind is created for them.
  \item[A6. Nothing applies.] At step 2 the student has no plan. The notification still announces the date and links the
        planner instead of listing courses.
\end{altflows}

\ucpart{Postconditions}
The notification is stored for the student; reading it marks it read. Nothing else changes.

\ucpart{Special requirements}
Notifications are in-app only in this release: there is no mail sender before B3. Dates come from the published school
calendar entered by the team. A notification never states a rule's value on its own; it shows the rule engine's result.

\ucpart{Relationships}
Includes UC-19. Links to UC-05 and UC-10.

%% ------------------------------------------------------------------ UC-13
\ucspec{UC-13}{Sign in}{F1 Account and onboarding}{A. Account and access}{Visitor}{Google identity (A2)}{High}{B0}{Phan
Bình Nguyên Lâm}

\ucpart{Brief description}
A visitor with an account signs in on a device and becomes a signed-in student. Signing in also cancels a pending account
deletion.

\ucpart{Preconditions}
The visitor has an account and is not signed in on this device.

\ucpart{Basic flow}
\begin{enumerate}
  \item The visitor chooses \st{Sign in}.
  \item The system asks for email and password.
  \item The visitor enters them and presses \st{Sign in}.
  \item The system checks them against the stored hash, creates a session for this device and signs the visitor in.
  \item If onboarding is finished, the system opens the dashboard (UC-03); otherwise it resumes onboarding (UC-01 step 6).
\end{enumerate}

\ucpart{Alternative flows}
\begin{altflows}
  \item[A1. Wrong email or password.] At step 4 the email is unknown or the password is wrong. The system says only that
        the email or password is invalid, without saying which, and answers in the same time in both cases.
  \item[A2. Sign in with Google] (after B0). At step 1 the visitor chooses \st{Continue with Google}. Google confirms the
        identity and the flow continues at step 5.
  \item[A3. Too many attempts.] At step 3 one address sends more than 10 sign-in requests in a minute. The system refuses
        more attempts until the minute has passed and says so.
  \item[A4. Account scheduled for deletion.] At step 4 the account was scheduled for deletion less than 30 days ago. The
        system cancels the deletion, signs the student in, and says the account has been restored. Shared links stay
        stopped (UC-09 A5).
\end{altflows}

\ucpart{Postconditions}
\textbf{Success:} a session exists for this device, valid for 7 days, rotated as it is used. \textbf{Failure:} no
session.

\ucpart{Special requirements}
Rate limiting per address. No hint of which field was wrong. Sign-in pages work signed out and by keyboard.

\ucpart{Relationships}
Leads to UC-03 or back into UC-01.

%% ------------------------------------------------------------------ UC-14
\ucspec{UC-14}{Manage account and privacy}{F1 Account and onboarding; privacy NFR}{A. Account and access}{Student}{none}%
{High}{B1}{Phan Bình Nguyên Lâm}

\ucpart{Brief description}
The signed-in student controls their own account: they sign out, see and revoke the devices signed in, change their
onboarding answers and preferences, withdraw AI consent, export everything OmniCourse holds about them, or delete the
account, with 30 days to change their mind.

\ucpart{Preconditions}
The student is signed in.

\ucpart{Basic flow}
\begin{enumerate}
  \item The student opens \st{Settings}.
  \item The system shows the sections: profile and planning goal, sessions, AI consent, appearance and language, export,
        and delete account.
  \item The student opens \st{Sessions}.
  \item The system lists every device signed in, with the browser, the place it was last used and the time.
  \item The student presses \st{Sign out} beside a device they do not recognise.
  \item The system revokes that session at once; that device must sign in again.
\end{enumerate}

\ucpart{Alternative flows}
\begin{altflows}
  \item[A1. Sign out here.] The student presses \st{Sign out} in the sidebar. The system ends this device's session and
        opens the landing page.
  \item[A2. Change program, cohort or goal.] At step 3 the student edits the profile and planning goal. The system saves
        them; existing plans keep the curriculum version they were made with, and the audit follows the new choice.
  \item[A3. Withdraw AI consent.] At step 3 the student turns off AI consent. The system stops sending their record at
        once; the advisor chat continues in general mode (UC-08 A1).
  \item[A4. Export the account.] At step 3 the student presses \st{Export my data}. The system builds one file with every
        row that belongs to them (account, transcript, plans, notifications, consent times) and no credential column, and
        the browser saves it.
  \item[A5. Delete the account.] At step 3 the student presses \st{Delete account}. The system explains the 30-day grace
        period and asks for the password again. On confirmation it signs out every device, stops every shared link and
        schedules the purge for 30 days later. Signing in during that time cancels it (UC-13 A4).
  \item[A6. Theme and language.] At step 3 the student switches light or dark theme, or Vietnamese or English. The change
        applies at once and is remembered on this device.
\end{altflows}

\ucpart{Postconditions}
The change asked for is applied; nothing else changes. After a completed deletion, every row with the student's id is
gone.

\ucpart{Special requirements}
Export and deletion cover every table that references the user (ADR 0004); a test fails on any table that would escape
them. The deletion needs the password again.

\ucpart{Relationships}
Changes what UC-08 may send; deletion affects UC-09 and UC-16.

%% ------------------------------------------------------------------ UC-15
\ucspec{UC-15}{View grades and what each course still needs}{F3 Degree audit dashboard}{B. Academic record}{Student}{none}%
{High}{B1}{Võ Nguyễn Minh Triết}

\ucpart{Brief description}
The student reads their grades term by term, with what each graded part of a course contributes, and, for a course not
finished yet, the minimum still needed to pass and to reach excellent. They can filter by requirement group and search.

\ucpart{Preconditions}
The student is signed in and has a transcript (UC-02).

\ucpart{Basic flow}
\begin{enumerate}
  \item The student opens \st{Grades}.
  \item The system shows three figures (cumulative GPA, credits passed and best term) beside a line chart of the term
        averages, with the pass mark (5) and the excellent mark (8) drawn as lines. A term with a failed course has a red
        point.
  \item Below, the system lists the current term, then the planned terms, then the past terms, newest first. Each course
        shows a bar from 0 to 10 with what each graded part adds stacked on it, thin marks at 5 and 8, the grade and a
        status: in progress, planned, passed or failed.
  \item For a course in progress, the system asks the rule engine (UC-19) what the parts still to be graded need, and shows,
        for example, \emph{Final exam needs 1.4 to pass, 7.4 for excellent}.
  \item The student opens a course.
  \item The system shows each graded part with its weight and score and, for a retaken course, every sitting.
\end{enumerate}

\ucpart{Alternative flows}
\begin{altflows}
  \item[A1. Filter by group.] At step 3 the student picks one or more requirement groups. The list keeps only those courses
        and hides empty terms.
  \item[A2. Search.] At step 3 the student types a code or a name. The list narrows as they type; if nothing matches, the
        system says so.
  \item[A3. Jump to a term.] At step 2 the student clicks a point of the chart. The page scrolls to that term.
  \item[A4. Already passed, or out of reach.] At step 4 the parts already graded secure a pass, or no score on the
        remaining parts can reach it. The system says \emph{Already sure to pass} or \emph{Cannot pass this sitting}
        instead of a minimum.
  \item[A5. Part scores unknown.] At step 4 the student has entered no part scores. The system shows the average still
        needed over the whole course: 5.0 to pass and 8.0 for excellent.
\end{altflows}

\ucpart{Postconditions}
Nothing is changed.

\ucpart{Special requirements}
Figures are in tabular numerals and read heavier than the text around them; units are written in full (\emph{16 credits},
not \emph{16 cr}). The minimum needed is computed by the rule engine from the course's assessment weights.

\ucpart{Relationships}
Extends UC-03. Includes UC-19.

%% ------------------------------------------------------------------ UC-16
\ucspec{UC-16}{View a shared plan}{F9 Plan export and sharing}{C. Planning}{Advisor}{none}{Low}{B2}{Nguyễn Anh Dũng}

\ucpart{Brief description}
An advisor opens the link a student sent and reads the plan, read-only, in the same layout the student uses, with the
violations the rule engine finds. No account is needed.

\ucpart{Preconditions}
The student has created a live link for the plan (UC-09).

\ucpart{Basic flow}
\begin{enumerate}
  \item The advisor opens the link.
  \item The system finds the live plan for the link and asks the rule engine (UC-19) to evaluate it against the curriculum
        version it was made with.
  \item The system shows a banner, \emph{Shared plan, read-only, last changed 30/10/2026 at 21:14}, then the plan: name,
        program, cohort and curriculum version; a status strip with the estimated graduation term or the number of
        violations, and the coverage of every requirement group; one row per term with its courses and credit meter; and
        a \st{Violations} panel.
  \item The advisor reviews the plan and replies to the student outside OmniCourse.
\end{enumerate}

\ucpart{Alternative flows}
\begin{altflows}
  \item[A1. Dead or unknown link.] At step 2 the link does not match a live shared plan: it was stopped, the plan or
        account was deleted, or it was mistyped. The system shows \emph{This plan is not shared}, the same page in every
        case, so it reveals nothing about whether a plan existed.
  \item[A2. Print.] At step 4 the advisor presses \st{Export PDF}. The browser prints the read-only view.
  \item[A3. Plan has violations.] At step 3 the plan breaks rules. The status strip says \emph{Not valid} with the count,
        and each violation shows its rule and clause.
\end{altflows}

\ucpart{Postconditions}
Nothing is changed. The page shows no grades, email, student ID, transcript or other plans.

\ucpart{Special requirements}
Public route under ADR 0010: read only, rate-limited (120 requests a minute per address), \texttt{noindex}. The view
follows the plan as it is now, not a copy made when the link was created.

\ucpart{Relationships}
Includes UC-19. Reached from UC-09.

%% ------------------------------------------------------------------ UC-17
\ucspec{UC-17}{Maintain the course catalog}{F12 Curriculum administration}{E. Curriculum administration}{Catalog
maintainer}{none}{Medium}{B2}{Nguyễn Hữu Quốc Thắng}

\ucpart{Brief description}
A catalog maintainer creates shared course records and corrects their name, credits or time breakdown. A course record is
shared by every program and cohort that links it, so a correction is shown with its full impact before it is saved.

\ucpart{Preconditions}
The user is signed in and holds catalog-maintainer permission (UC-18).

\ucpart{Basic flow}
\begin{enumerate}
  \item The catalog maintainer opens \st{Course catalog}. No program needs to be selected.
  \item The system lists the courses with code, name and credits, and a search box.
  \item The maintainer opens an existing course and corrects, for example, its credits.
  \item The system shows the old and new values and every program and cohort that links the course, and explains that all
        of them will read the new values.
  \item The maintainer checks the official source and presses \st{Confirm and save}.
  \item The system checks the permission, re-checks the list of linked curricula, and saves.
\end{enumerate}

\ucpart{Alternative flows}
\begin{altflows}
  \item[A1. Create a course.] At step 3 the maintainer presses \st{New course} and enters a code, name and credits. The
        system checks the code format and rejects a duplicate code.
  \item[A2. Impact changed during review.] At step 6 another curriculum linked the course since step 4. The system shows
        the new impact and asks for confirmation again.
  \item[A3. Value differs by cohort.] At step 3 the correction applies to some cohorts only. The shared record cannot
        hold that; the system refuses and says a course version is needed (a design still open).
  \item[A4. Delete.] At step 3 the maintainer deletes a course. Only an unlinked course may be deleted; a linked one is
        refused with the list of links.
  \item[A5. Program permission only.] A program administrator without catalog permission tries to save. The server
        refuses (UC-11 A4).
\end{altflows}

\ucpart{Postconditions}
Every linked program and cohort reads the new values; the audit and planner may give different results from then on.
Stored transcripts and plans are not rewritten.

\ucpart{Special requirements}
The impact list is computed on the server at save time, not trusted from the screen.

\ucpart{Relationships}
Extends UC-11 at its step 3 (A8).

%% ------------------------------------------------------------------ UC-18
\ucspec{UC-18}{Manage administration access}{F12 Curriculum administration}{E. Curriculum administration}{Access
manager}{none}{Medium}{B2}{Nguyễn Hữu Quốc Thắng}

\ucpart{Brief description}
An access manager creates programs and grants or revokes the three administration permissions: program administrator for
a named program, catalog maintainer, and access manager. Permissions are separate and denied by default.

\ucpart{Preconditions}
The user is signed in and holds access-manager permission. The first access manager is set up by the Project Lead through a
controlled setup step.

\ucpart{Basic flow}
\begin{enumerate}
  \item The access manager opens \st{Administration access}.
  \item The system lists the users who hold a permission, with the permission and, for program administrators, the
        program.
  \item The access manager presses \st{Grant}, finds an existing user by email, and picks a permission and, if needed, a
        program.
  \item The system checks that the user exists and that the grant is new, saves it and shows it in the list.
\end{enumerate}

\ucpart{Alternative flows}
\begin{altflows}
  \item[A1. Revoke.] At step 3 the access manager presses \st{Revoke} on a grant and confirms. The system removes it; the
        next write by that user is refused.
  \item[A2. Add a program.] At step 3 the access manager presses \st{Add program} and enters its code and name. The system
        rejects a duplicate code. Creating a program grants nobody the right to edit it.
  \item[A3. Unknown user.] At step 4 no account has that email. The system says the person must sign up first; it never
        creates an account here.
  \item[A4. Revoke your own last access-manager grant.] The system refuses, so the system always keeps an access manager.
\end{altflows}

\ucpart{Postconditions}
The grant or revocation applies to the next request the user sends. Sign-up and onboarding never grant a permission.

\ucpart{Special requirements}
Every write in UC-11 and UC-17 checks the current grants on the server, including just after a revocation.

\ucpart{Relationships}
Gives the rights that UC-11 and UC-17 require.

%% ------------------------------------------------------------------ UC-19
\ucspec{UC-19}{Evaluate academic rules}{F3, F4, F5, F6, F7 (shared)}{All modules (included)}{none: included by other use
cases}{none}{High}{B1 audit; B2 plan checks}{Võ Nguyễn Minh Triết}

\ucpart{Brief description}
The rule engine decides every academic rule for a student: which transcript attempts count and pass, how credits fill each
requirement group, whether a course's prerequisites are met in time, whether a term offers the course, whether a term's load
is within the credit cap, what an unfinished course still needs, and when the student can graduate. It is included by every
use case that shows or uses such a result, so all of them agree.

\ucpart{Preconditions}
The curriculum version of the student's cohort is available. A transcript and a plan are optional inputs.

\ucpart{Basic flow}
\begin{enumerate}
  \item The including use case asks for an evaluation and passes the cohort's curriculum, the student's transcript rows
        and, if any, a plan.
  \item The rule engine picks, for each course, the attempt that counts under the regulation, and marks each attempt passed
        or failed.
  \item It assigns passed and in-progress credits to the requirement groups and computes, per group, passed, in progress
        and still needed.
  \item For each course placed in the plan, it checks that every prerequisite is passed or placed in an earlier term, and
        that the term offers the course.
  \item For each term, it compares the credit load with the program's minimum and cap for that kind of term.
  \item It computes the cumulative GPA and the earliest term in which every requirement is met: the estimated graduation
        term.
  \item It returns the results and a list of violations. Each violation has a code, the courses or term involved, a
        sentence, and the regulation clause it comes from.
\end{enumerate}

\ucpart{Alternative flows}
\begin{altflows}
  \item[A1. Check one change.] At step 1 the caller passes a plan and one proposed move (UC-05 step 4). The engine
        evaluates the plan as it would be after the move and returns the differences: new and resolved violations, and the
        new load of each affected term.
  \item[A2. What a course still needs.] At step 1 the caller passes one course in progress with its graded parts (UC-15).
        The engine returns the minimum score on the remaining parts for a pass and for excellent, or that the result is
        already sure or out of reach.
  \item[A3. Data missing.] At any step a rule needs data the curriculum does not have, such as offering terms. The engine
        returns \emph{unknown} for that check, never a guess, and the caller says so.
  \item[A4. Curriculum check.] At step 1 the caller is UC-11 and passes a changed curriculum. The engine checks it for a
        prerequisite cycle and for groups whose required credits cannot be reached.
\end{altflows}

\ucpart{Postconditions}
Nothing is stored. The same inputs always give the same results.

\ucpart{Special requirements}
\begin{itemize}
  \item Deterministic code, no language model (hard rule 1). A model may call the engine and explain its results, never
        override them.
  \item At least 80\% statement coverage, enforced in CI; tests per requirement group with synthetic transcripts.
  \item A plan change is evaluated in under 300 ms.
\end{itemize}

\ucpart{Relationships}
Included by UC-03, UC-04, UC-05, UC-06, UC-07, UC-08, UC-11, UC-12, UC-15 and UC-16.
